\chapter*{Wykaz ważniejszych oznaczeń i skrótów}
\addcontentsline{toc}{chapter}{Wykaz ważniejszych oznaczeń i skrótów}

\begin{itemize}
  \item \textbf{API} (ang. \textit{Application Programming Interface}) - zestaw reguł i protokołów określających sposób komunikacji oraz integracji pomiędzy różnymi elementami aplikacji.
  \item \textbf{AR} (ang. \textit{Augmented Reality}) - technologia rozszerzonej rzeczywistości, która łączy elementy wirtualne z rzeczywistym światem użytkownika.
  \item \textbf{CI/CD} (ang. \textit{Continuous Integration / Continuous Delivery lub Deployment}) - zestaw praktyk automatyzujących budowanie, testowanie i dostarczanie oprogramowania, dzięki czemu zmiany trafiają do użytkowników szybciej i bardziej niezawodnie.
  \item \textbf{CORS} (ang. \textit{Cross-Origin Resource Sharing}) - mechanizm kontrolujący, które domeny mogą wykonywać zapytania do zasobów serwera.
  \item \textbf{CQRS} (ang. \textit{Command Query Responsibility Segregation}) - wzorzec architektoniczny polegający na rozdzieleniu operacji modyfikujących stan systemu od operacji odczytujących dane.
  \item \textbf{DDD} (ang. \textit{Domain-Driven Design}) - podejście do projektowania oprogramowania, które skupia sie na odwzorowywaniu rzeczywistego świata w modelu domeny.
  \item \textbf{DOM} (ang. \textit{Document Object Model}) - struktura reprezentująca dokumenty HTML i XML jako hierarchię obiektów, umożliwiająca manipulację ich zawartością i strukturą za pomocą języków programowania.
  \item \textbf{DTO} (ang. \textit{Data Transfer Object}) - obiekt służący do przenoszenia danych między warstwami aplikacji.
  \item \textbf{E2E} (ang. \textit{End-to-End}) - rodzaj testów sprawdzających cały przepływ działania aplikacji z perspektywy użytkownika.
  \item \textbf{ERD} (ang. \textit{Entity-Relationship Diagram}) - diagram związków encji, służący do modelowania struktury bazy danych.
  \item \textbf{GPS} (ang. \textit{Global Positioning System}) - system nawigacji satelitarnej umożliwiający określenie położenia geograficznego na powierzchni Ziemi.
  \item \textbf{GUI} (ang. \textit{Graphical User Interface}) - graficzny interfejs użytkownika oparty na elementach wizualnych, takich jak okna, ikony czy przyciski.
  \item \textbf{HTML} (ang. \textit{HyperText Markup Language}) - język znaczników używany do tworzenia stron internetowych.
  \item \textbf{HSTS} (ang. \textit{HTTP Strict Transport Security}) - mechanizm wymuszający użycie HTTPS, chroniący przed atakami typu downgrading i przechwytywaniem sesji.
  \item \textbf{HTTP} (ang. \textit{HyperText Transfer Protocol}) - podstawowy protokół komunikacji w sieci WWW, umożliwiający przesyłanie stron internetowych oraz danych między klientem a serwerem.
  \item \textbf{HTTPS} (ang. \textit{HyperText Transfer Protocol Secure}) - bezpieczna wersja protokołu HTTP, która wykorzystuje szyfrowanie SSL/TLS do ochrony przesyłanych danych.
  \item \textbf{IDE} (ang. \textit{Integrated Development Environment}) - zintegrowane środowisko programistyczne zawierające narzędzia potrzebne do tworzenia i debugowania kodu.
  \item \textbf{ITS} (ang. \textit{Issue Tracker System}) - sposób zarządzania systemem, który odpowiada na zapytania wysyłane dowolną drogą.
  \item \textbf{JSON} (ang. \textit{JavaScript Object Notation}) - format wymiany danych oparty na składni języka JavaScript.
  \item \textbf{JWT} (ang. \textit{JSON Web Token}) - standardowy format bezpiecznego przesyłania informacji jako obiekt JSON.
  \item \textbf{LINQ} (ang. \textit{Language Integrated Query}) - zestaw technik umożliwiających tworzenie zapytań bezpośrednio w kodzie C\#, w spójny sposób niezależnie od źródła danych.
  \item \textbf{NPM} (ang. \textit{Node Package Manager}) - manager pakietów dedykowany środowisku Node.js.
  \item \textbf{ORM} (ang. \textit{Object-Relational Mapping}) - technika mapowania danych między systemem obiektowym a relacyjną bazą danych.
  \item \textbf{PWA} (ang. \textit{Progressive Web Application}) - aplikacja webowa wykorzystująca nowoczesne technologie webowe w celu zapewnienia doświadczenia podobnego do aplikacji natywnych.
  \item \textbf{QA} (ang. \textit{Quality Assurance}) - zestaw procesów i praktyk mających zapewnić wysoką jakość tworzonego oprogramowania.
  \item \textbf{REST} (ang. \textit{Representational State Transfer}) - styl architektoniczny wykorzystywany przy projektowaniu usług sieciowych.
  \item \textbf{SQL} (ang. \textit{Structured Query Language}) - strukturalny język zapytań służący do definiowania, modyfikowania i pobierania danych z relacyjnych baz danych.
  \item \textbf{TLS} (ang. \textit{Transport Layer Security}) - protokół szyfrujący zapewniający poufność i integralność danych przesyłanych w sieci.
  \item \textbf{UI} (ang. \textit{User Interface}) - interfejs użytkownika, czyli sposób prezentacji i organizacji elementów umożliwiających interakcję z aplikacją.
  \item \textbf{UX} (ang. \textit{User Experience}) - doświadczenie użytkownika obejmujące wrażenia, odczucia i wygodę korzystania z produktu lub systemu.
  \item \textbf{WWW} (ang. \textit{World Wide Web}) - system zasobów internetowych połączonych hiperłączami, umożliwiający przeglądanie stron i treści za pomocą przeglądarki.
  \item \textbf{VCS} (ang. \textit{Version Control System}) - system kontroli wersji pozwalający śledzić zmiany w kodzie, zarządzać historią projektu i pracą zespołową.
  \item \textbf{VM} (ang. \textit{Virtual Machine}) - maszyna wirtualna emulująca komputer w celu uruchamiania systemów operacyjnych lub aplikacji w izolowanym środowisku.
  \item \textbf{VS} (ang. \textit{Visual Studio}) - zintegrowane środowisko programistyczne (IDE) stworzone przez firmę Microsoft
  \item \textbf{YAML} (ang. \textit{YAML Ain't Markup Language}) - lekki format plików konfiguracyjnych oparty na strukturze wcięć.
\end{itemize}
